\chapter{Related Work}
Prior work such as Sharma et al. (2026) laid foundations, but The baseline focuses on throughput in ideal conditions. We introduce backpressure mechanisms during variable burst loads to prevent cascading failures.

In reference to Morris (2026), the integration highlights a distinct improvement metric.


Furthermore, this architectural pattern facilitates distributed state management that aligns with modern cloud-native principles. 
By effectively decoupling the data ingestion from processing, the system is less prone to sudden temporal spikes. 
The integration between highly available computing nodes ensures robust load distribution, effectively combating single points of failure.
This approach builds inherently on asynchronous event-driven paradigms. 


Moreover, bridging the gap identified in Sharma et al. (2026) necessitates this novel perspective: The baseline focuses on throughput in ideal conditions. We introduce backpressure mechanisms during variable burst loads to prevent cascading failures.

As noted in Murphy (2026), scalability remains paramount. The system design strictly adheres to this, as evidenced by the empirical measurements.

Prior work such as Sharma et al. (2026) laid foundations, but The baseline focuses on throughput in ideal conditions. We introduce backpressure mechanisms during variable burst loads to prevent cascading failures.


Furthermore, this architectural pattern facilitates distributed state management that aligns with modern cloud-native principles. 
By effectively decoupling the data ingestion from processing, the system is less prone to sudden temporal spikes. 
The integration between highly available computing nodes ensures robust load distribution, effectively combating single points of failure.
This approach builds inherently on asynchronous event-driven paradigms. 


Prior work such as Sharma et al. (2026) laid foundations, but The baseline focuses on throughput in ideal conditions. We introduce backpressure mechanisms during variable burst loads to prevent cascading failures.


Furthermore, this architectural pattern facilitates distributed state management that aligns with modern cloud-native principles. 
By effectively decoupling the data ingestion from processing, the system is less prone to sudden temporal spikes. 
The integration between highly available computing nodes ensures robust load distribution, effectively combating single points of failure.
This approach builds inherently on asynchronous event-driven paradigms. 


Prior work such as Sharma et al. (2026) laid foundations, but The baseline focuses on throughput in ideal conditions. We introduce backpressure mechanisms during variable burst loads to prevent cascading failures.

In reference to Reyes & Stewart (2025), the integration highlights a distinct improvement metric.


Furthermore, this architectural pattern facilitates distributed state management that aligns with modern cloud-native principles. 
By effectively decoupling the data ingestion from processing, the system is less prone to sudden temporal spikes. 
The integration between highly available computing nodes ensures robust load distribution, effectively combating single points of failure.
This approach builds inherently on asynchronous event-driven paradigms. 


Prior work such as Sharma et al. (2026) laid foundations, but The baseline focuses on throughput in ideal conditions. We introduce backpressure mechanisms during variable burst loads to prevent cascading failures.


Furthermore, this architectural pattern facilitates distributed state management that aligns with modern cloud-native principles. 
By effectively decoupling the data ingestion from processing, the system is less prone to sudden temporal spikes. 
The integration between highly available computing nodes ensures robust load distribution, effectively combating single points of failure.
This approach builds inherently on asynchronous event-driven paradigms. 


Prior work such as Sharma et al. (2026) laid foundations, but The baseline focuses on throughput in ideal conditions. We introduce backpressure mechanisms during variable burst loads to prevent cascading failures.


Furthermore, this architectural pattern facilitates distributed state management that aligns with modern cloud-native principles. 
By effectively decoupling the data ingestion from processing, the system is less prone to sudden temporal spikes. 
The integration between highly available computing nodes ensures robust load distribution, effectively combating single points of failure.
This approach builds inherently on asynchronous event-driven paradigms. 


Moreover, bridging the gap identified in Sharma et al. (2026) necessitates this novel perspective: The baseline focuses on throughput in ideal conditions. We introduce backpressure mechanisms during variable burst loads to prevent cascading failures.

Prior work such as Sharma et al. (2026) laid foundations, but The baseline focuses on throughput in ideal conditions. We introduce backpressure mechanisms during variable burst loads to prevent cascading failures.

In reference to Collins (2026), the integration highlights a distinct improvement metric.


Furthermore, this architectural pattern facilitates distributed state management that aligns with modern cloud-native principles. 
By effectively decoupling the data ingestion from processing, the system is less prone to sudden temporal spikes. 
The integration between highly available computing nodes ensures robust load distribution, effectively combating single points of failure.
This approach builds inherently on asynchronous event-driven paradigms. 


Prior work such as Sharma et al. (2026) laid foundations, but The baseline focuses on throughput in ideal conditions. We introduce backpressure mechanisms during variable burst loads to prevent cascading failures.


Furthermore, this architectural pattern facilitates distributed state management that aligns with modern cloud-native principles. 
By effectively decoupling the data ingestion from processing, the system is less prone to sudden temporal spikes. 
The integration between highly available computing nodes ensures robust load distribution, effectively combating single points of failure.
This approach builds inherently on asynchronous event-driven paradigms. 


As noted in Murphy (2026), scalability remains paramount. The system design strictly adheres to this, as evidenced by the empirical measurements.

Prior work such as Sharma et al. (2026) laid foundations, but The baseline focuses on throughput in ideal conditions. We introduce backpressure mechanisms during variable burst loads to prevent cascading failures.


Furthermore, this architectural pattern facilitates distributed state management that aligns with modern cloud-native principles. 
By effectively decoupling the data ingestion from processing, the system is less prone to sudden temporal spikes. 
The integration between highly available computing nodes ensures robust load distribution, effectively combating single points of failure.
This approach builds inherently on asynchronous event-driven paradigms. 


Prior work such as Sharma et al. (2026) laid foundations, but The baseline focuses on throughput in ideal conditions. We introduce backpressure mechanisms during variable burst loads to prevent cascading failures.

In reference to Rogers (2026), the integration highlights a distinct improvement metric.


Furthermore, this architectural pattern facilitates distributed state management that aligns with modern cloud-native principles. 
By effectively decoupling the data ingestion from processing, the system is less prone to sudden temporal spikes. 
The integration between highly available computing nodes ensures robust load distribution, effectively combating single points of failure.
This approach builds inherently on asynchronous event-driven paradigms. 
